% \iffalse
%<*metadata>
% \fi
% \section{Metadata}
% \subsection{Private Storage Macros}
%    \begin{macrocode}
\makeatletter
\newcommand{\@authorid}{}
\newcommand{\@firstadvisor}{}
\newcommand{\@firstadvisorid}{}
\newcommand{\@secondadvisor}{}
\newcommand{\@secondadvisorid}{}
\newcommand{\@firstexaminer}{}
\newcommand{\@firstexaminerid}{}
\newcommand{\@secondexaminer}{}
\newcommand{\@secondexaminerid}{}
\newcommand{\@headdepartment}{}
\newcommand{\@headdepartmentid}{}
\newcommand{\@faculty}{}
\newcommand{\@program}{}
\newcommand{\@university}{}
\newcommand{\@city}{}
\newcommand{\@yearofsubmission}{}
%    \end{macrocode}
%    
% \subsection{Setter Macros}
% Define user commands with optional ID arguments to set these values.
%    \begin{macrocode}
\newcommand{\firstadvisor}[2][]{\gdef\@firstadvisor{#2}\gdef\@firstadvisorid{#1}}
\newcommand{\secondadvisor}[2][]{\gdef\@secondadvisor{#2}\gdef\@secondadvisorid{#1}}
\newcommand{\firstexaminer}[2][]{\gdef\@firstexaminer{#2}\gdef\@firstexaminerid{#1}}
\newcommand{\secondexaminer}[2][]{\gdef\@secondexaminer{#2}\gdef\@secondexaminerid{#1}}
\newcommand{\headdepartment}[2][]{\gdef\@headdepartment{#2}\gdef\@headdepartmentid{#1}}
%    \end{macrocode}
% Commands without ID arguments
%    \begin{macrocode}
\newcommand{\faculty}[1]{\gdef\@faculty{#1}}
\newcommand{\program}[1]{\gdef\@program{#1}}
\newcommand{\university}[1]{\gdef\@university{#1}}
\newcommand{\city}[1]{\gdef\@city{#1}}
\newcommand{\yearofsubmission}[1]{\gdef\@yearofsubmission{#1}}
%    \end{macrocode}
%    
% \subsection{Getter Macros}
% Now expose robust “the‑” getters to access the values
%    \begin{macrocode}
\newcommand{\theauthorid}{\@authorid}
\newcommand{\thefirstadvisor}{\@firstadvisor}
\newcommand{\thefirstadvisorid}{\@firstadvisorid}
\newcommand{\thesecondadvisor}{\@secondadvisor}
\newcommand{\thesecondadvisorid}{\@secondadvisorid}
\newcommand{\thefirstexaminer}{\@firstexaminer}
\newcommand{\thefirstexaminerid}{\@firstexaminerid}
\newcommand{\thesecondexaminer}{\@secondexaminer}
\newcommand{\thesecondexaminerid}{\@secondexaminerid}
\newcommand{\theheaddepartment}{\@headdepartment}
\newcommand{\theheaddepartmentid}{\@headdepartmentid}
\newcommand{\thefaculty}{\@faculty}
\newcommand{\theprogram}{\@program}
\newcommand{\theuniversity}{\@university}
\newcommand{\thecity}{\@city}
\newcommand{\theyearofsubmission}{\@yearofsubmission}
%    \end{macrocode}
% \subsection{Metadata Validation}
% \begin{macro}{\@check}
% Define internal macro to check if required metadata fields are set. 
%    \begin{macrocode}
\newcommand{\@check}[1]{%
%    \end{macrocode}
% The extra |\expandafter| forces the |\csname| to be fully expanded before |\ifdefstring| sees it.
%    \begin{macrocode}
  \expandafter\ifdefstring\expandafter{\csname @#1\endcsname}{}{%
%    \end{macrocode}
% If the field is empty, raise a class error prompting the user to define it in the preamble.
%    \begin{macrocode}
    \ClassError{skripsi.cls}{Missing \protect#1}%
      {You must declare \protect#1{...} in preamble}%
  }%
}
%    \end{macrocode}
% \end{macro}
% 
% \subsection{Title and Author}
% Redefining primitive title and author commands is implemented quite different to keep the compatibility with titling package
%    \begin{macrocode}
% Load required packages
\RequirePackage{letltxmacro}
%    \end{macrocode}
% by loading |letltxmacro| to save the titling-modified \title command
%    \begin{macrocode}
\LetLtxMacro\titlingoriginal\title
\LetLtxMacro\authororiginal\author
%    \end{macrocode}
% This happens AFTER |titling| package is loaded. Then define keys for title language variants
%    \begin{macrocode}
\define@key{titlekeys}{en}{\gdef\@title@en{#1}}
\define@key{titlekeys}{id}{\gdef\@title@id{#1}}
\define@key{titlekeys}{fr}{\gdef\@title@fr{#1}}
\define@key{titlekeys}{es}{\gdef\@title@es{#1}}
%    \end{macrocode}
% Initialize to empty
%    \begin{macrocode}
\gdef\@title@en{}
\gdef\@title@id{}
\gdef\@title@fr{}
\gdef\@title@es{}
\gdef\@authorid{}
%    \end{macrocode}
% Then begin redefine |\title| to accept optional key-value arguments as the other language title variants.
%    \begin{macrocode}
\renewcommand{\title}[2][]{%
%    \end{macrocode}
% Process the optional key-value
%    \begin{macrocode}
  \setkeys{titlekeys}{#1}%
%    \end{macrocode}
% Call titling's |\title| (defines |\thetitle|), so |\thetitle|
%    \begin{macrocode}
  \titlingoriginal{#2}%        
}
%    \end{macrocode}
% also redefine |\author| to accept
%    \begin{macrocode}
\renewcommand{\author}[2][]{%
%    \end{macrocode}
% optional arguments to store author ID.
%    \begin{macrocode}
  \gdef\@authorid{#1}%
%    \end{macrocode}
% while keeping the original functionality of titling's |\author| that can be accessed through |\theauthor|.
%    \begin{macrocode}
  \authororiginal{#2}%
}
%    \end{macrocode}
%    
% \subsection{Getter Macros for Title and Author}
% User-friendly commands to access the titles through |\thetitleXX| macros where XX is two letters language code.
%    \begin{macrocode}
\newcommand{\thetitleen}{\@title@en}
\newcommand{\thetitleid}{\@title@id}
\newcommand{\thetitlefr}{\@title@fr}
\newcommand{\thetitlees}{\@title@es}
\makeatother
%    \end{macrocode}
% \iffalse
%</metadata>
% \fi